\documentclass[11pt]{article}
\usepackage[a4paper,margin=25mm]{geometry}
\usepackage[T1]{fontenc}
\usepackage{lmodern}
\usepackage{amsmath,amssymb,amsthm,mathtools}
\usepackage{booktabs,longtable}
\usepackage[hidelinks]{hyperref}
\usepackage{microtype}
\hypersetup{pdftitle={Schur leading terms and a toric degeneration of the Lecture Hall algebra},pdfauthor={Aurantium},pdfsubject={Lecture Hall initial algebra, explicit subset map, and integral flat family}}
\newtheorem{theorem}{Theorem}[section]
\newtheorem{lemma}[theorem]{Lemma}
\newtheorem{proposition}[theorem]{Proposition}
\newtheorem{corollary}[theorem]{Corollary}
\theoremstyle{definition}
\newtheorem{definition}[theorem]{Definition}
\theoremstyle{remark}
\newtheorem{remark}[theorem]{Remark}
\newcommand{\Z}{\mathbb Z}
\newcommand{\Q}{\mathbb Q}
\newcommand{\N}{\mathbb Z_{\ge0}}
\newcommand{\LT}{\operatorname{LT}}
\newcommand{\inA}{\operatorname{in}}
\newcommand{\supp}{\operatorname{supp}}
\newcommand{\vol}{\operatorname{vol}}
\newcommand{\diag}{\operatorname{diag}}
\newcommand{\transpose}{\mathsf T}
\title{Schur leading terms and a toric degeneration\\of the Lecture Hall algebra}
\author{Aurantium}
\date{October 13, 2026}
\begin{document}
\maketitle
\begin{abstract}
Starting from the signed-moment identification of the Lecture Hall polynomials
with signed Stieltjes--Rogers moments, we determine the leading term of every
proposed SAGBI generator. The key identity is a highest-degree recursion for
Schur determinants of arbitrary Stieltjes--Rogers moments. Its iteration
reads off the Frobenius coordinates of the indexing partition. After the
Lecture Hall sign specialization, every leading coefficient is $1$, and the
resulting exponents give an explicit bijection from subsets of $[n-1]$ to
the Hilbert basis of the Lecture Hall cone.

We then prove that the initial algebra is the entire Lecture Hall semigroup
ring, for every $n$ and over every field. The last step uses equality of
elementary lattice-point growth constants, not the Lecture Hall Theorem.
Thus the argument proves Katth\"an's Conjecture 4.2 and answers Problem 5.1,
with the explicitly stated square-minor interpretation. It also produces a
monic integral SAGBI basis and an explicit free, hence flat, one-parameter
degeneration. All needed determinant and counting arguments are included.
Finite exact records and code accompany this note; they are not used as proofs.
No claim of publication priority is made.
\end{abstract}

\section{Definitions and statements}
All variables below are algebraically independent unless a specialization
is explicitly stated. Let $\mu_j(\alpha_1,\alpha_2,\ldots)$ be the sum of
weights of Dyck paths of semilength $j$: up steps have weight $1$ and a down
step from height $h$ has weight $\alpha_h$. Set $\mu_0=1$ and $\mu_j=0$ for
$j<0$. Define
\begin{equation}\label{eq:weights}
 \epsilon_h=\begin{cases}-1&h\ge4\text{ is even},\\1&\text{otherwise},\end{cases}
 \qquad \alpha_h=\epsilon_h y_hy_{h+1},\qquad
 \ell_{j+1}=y_1\mu_j(\alpha)\quad(j\ge0).
\end{equation}
Also set $\ell_j=0$ for $j\le0$.
The polynomiality note \cite{Polynomiality} establishes that these are the Lecture Hall polynomials.
Appendix~\ref{app:identification} reproves the needed identification so that
the present argument is self-contained.

For $S=\{s_1<\cdots<s_r\}\subseteq[n-1]$, $r>0$, put
\begin{equation}\label{eq:minor}
 \ell_S=(-1)^{r+1}\det[\ell_{s_j+2-i}]_{1\le i,j\le r},
 \qquad \ell_\varnothing=\ell_1=y_1.
\end{equation}
This is the negative of the top-row minor of the negative upper-triangular
Toeplitz matrix in Katth\"an's Definition 4.1 \cite{Katthan}.
Let
\[
 A_n(k)=k[\ell_1,\ldots,\ell_n]\subseteq k[y_1,\ldots,y_n],\qquad
 \mathcal B_n=\{\ell_S:S\subseteq[n-1]\}.
\]
The term order is degree-lexicographic, with $y_1>\cdots>y_n$.
For a polynomial its leading term includes its coefficient.
Write
\begin{equation}\label{eq:cone}
 C_n=\left\{x\in\mathbb R^n:
 \frac{x_1}{n}\ge\frac{x_2}{n-1}\ge\cdots\ge x_n\ge0\right\},
 \qquad L_n=C_n\cap\Z^n.
\end{equation}
The semigroup algebra $k[L_n]$ is the subalgebra spanned by the monomials
$y^e$ with $e\in L_n$.

\begin{theorem}[All dimensions; all characteristics]\label{thm:main}
For every $n\ge1$ and every field $k$,
\[
 \inA_{\mathrm{deglex}}A_n(k)=k[L_n].
\]
The set $\mathcal B_n$ is a minimal SAGBI basis. Its leading coefficients
are all $1$, and its leading exponents run bijectively through
\begin{equation}\label{eq:hilbert}
 \mathcal H_n=\{(1,0,\ldots,0)\}\ \cup\
 \{(a+1,a,b_3,\ldots,b_t,0,\ldots,0):
 1\le a\le n-1,\quad a>b_3>\cdots>b_t>0\},
\end{equation}
where the tail after $a$ may be empty. This set is the Hilbert basis of $L_n$.
\end{theorem}

Here is the promised explicit formula. For nonempty $S$ define a partition
of length $r$ by
\begin{equation}\label{eq:lambda}
 \lambda_i=s_{r+1-i}-r+i\qquad(1\le i\le r).
\end{equation}
Let $D$ be its Durfee size, and let
\[
 A_i=\lambda_i-i,\qquad B_i=\lambda'_i-i\quad(1\le i\le D)
\]
be its Frobenius arm and leg lengths. Form the sequence
\begin{equation}\label{eq:profile}
 (d_1,d_2,d_3,d_4,\ldots)=(A_1+1,B_1,A_2+1,B_2,\ldots,A_D+1,B_D),
\end{equation}
remove a final zero if present, and pad with zeros thereafter.

\begin{theorem}[Leading-term formula and the subset map]\label{thm:formula}
For every nonempty $S\subseteq[n-1]$,
\begin{equation}\label{eq:leading}
 \LT(\ell_S)=y_1^{r+d_1}\prod_{j=2}^n y_j^{d_{j-1}+d_j}.
\end{equation}
In particular the first two exponents are $\max(S)+1$ and $\max(S)$.
Together with $\LT(\ell_\varnothing)=y_1$, formula~\eqref{eq:leading}
is an explicit, invertible description of the map in Katth\"an's
Problem 5.1.
\end{theorem}

The source conventions matter. The square minor defining $E_i$ uses rows
$1,\ldots,\lceil i/2\rceil$ and columns $\lfloor i/2\rfloor+1,\ldots,i$.
This is the interpretation in \cite{Polynomiality}, supported by
the source's examples and Definition 4.1. We likewise retain the corrected
fourth polynomial consistent with the defining equations and the source's
Section 6. These are not represented as author-issued errata.

\section{A determinant identity from reciprocal series}
We first work with an arbitrary sequence $h_0=1,h_1,h_2,\ldots$ over a
commutative ring, with $h_j=0$ for $j<0$. Define $e_j$ by
\[
 \left(\sum_{j\ge0}h_jz^j\right)
 \left(\sum_{j\ge0}(-1)^je_jz^j\right)=1,
 \qquad e_0=1,
\]
and set $e_j=0$ for $j<0$.
For a partition $\lambda$ of length $r$, write
\[
 s_\lambda(h)=\det[h_{\lambda_i-i+j}]_{1\le i,j\le r},
 \qquad s_\varnothing(h)=1.
\]
Trailing zero parts do not change this determinant: the bottom row then
has a $1$ in its last position and zeros before it, and expansion removes
that row and column.

\begin{lemma}[Dual Jacobi--Trudi identity]\label{lem:dual}
For every partition $\lambda$,
\[
 s_\lambda(h)=s_{\lambda'}(e).
\]
\end{lemma}
\begin{proof}
The empty case is immediate. Set $L=\lambda_1$, $N=r+L$, and index the
$N\times N$ matrix $U$ by $0,\ldots,N-1$, with $U_{ab}=h_{a-b}$.
It is lower triangular with determinant $1$, and its inverse has entries
$(U^{-1})_{ab}=(-1)^{a-b}e_{a-b}$, by the reciprocal-series identity.

Use the row set $I=\{\lambda_i+r-i:1\le i\le r\}$ and column set
$J=\{0,\ldots,r-1\}$. Reversing both ordered index sets identifies the
minor $U_{I,J}$ with $s_\lambda(h)$. The complement of $I$ is, in increasing
order,
\[
 I^c=\{r+j-1-\lambda'_j:1\le j\le L\},\qquad
 J^c=\{r,\ldots,r+L-1\}.
\]
To check the first formula, its entries are strictly increasing, lie in
$[0,N-1]$, and cannot equal $\lambda_i+r-i$: equality would say
$\lambda_i+\lambda'_j=i+j-1$, whereas if $\lambda_i\ge j$ the left side
is at least $i+j$, and otherwise it is at most $i+j-2$.
The two sets therefore are complementary by cardinality.

The complementary-minor identity gives
\[
 \det U_{I,J}=(-1)^{\sum I+\sum J}\det(U^{-1})_{J^c,I^c}.
\]
Here $\sum I+\sum J\equiv |\lambda|\pmod2$.
The transposed complementary minor has entries
\[
 (-1)^{\lambda'_j-j+i}e_{\lambda'_j-j+i}.
\]
Their determinant has the overall sign $(-1)^{|\lambda|}$, so the two
signs cancel.

For completeness, the complementary-minor identity follows by permuting
selected rows and columns to the front and writing the resulting matrix
in four blocks $A,B,C,D$, in row order.
When $A$ is invertible, the determinant of the whole matrix is
$\det A\det(D-CA^{-1}B)$, while the corresponding bottom-right block of
the inverse is $(D-CA^{-1}B)^{-1}$. The permutation signs give the displayed
formula. Proving this first for a generic matrix and clearing denominators
makes it a polynomial identity, hence valid for the present matrix as well.
\end{proof}

Return now to the Dyck moments with independent weights $\alpha_h$.
Write $h_j=\mu_j(\alpha_1,\alpha_2,\ldots)$ and
$\widehat h_j=\mu_j(\alpha_2,\alpha_3,\ldots)$.
First-return decomposition of a Dyck path gives the formal identity
\begin{equation}\label{eq:reciprocal}
 H(z):=\sum_{j\ge0}h_jz^j
 =\frac{1}{1-\alpha_1z\widehat H(z)},\qquad
 \widehat H(z)=\sum_{j\ge0}\widehat h_jz^j.
\end{equation}
Consequently
\begin{equation}\label{eq:e}
 e_j=(-1)^{j-1}\alpha_1\widehat h_{j-1}\qquad(j\ge1).
\end{equation}
All these are identities of formal power series; a coefficient of order
$j$ uses only finitely many heights.

\begin{lemma}[Highest-degree recursion]\label{lem:peel}
For nonempty $\lambda$, set $L=\lambda_1$ and
$\nu=(\lambda_2,\lambda_3,\ldots)'$. Then
\begin{align}
 \deg_{\alpha_1}s_\lambda(\mu)&=L,\label{eq:degreealpha}\\
 [\alpha_1^L]s_\lambda(\mu)
 &=(-1)^{|\lambda|-L}s_\nu(\widehat\mu).\label{eq:peel}
\end{align}
\end{lemma}
\begin{proof}
By Lemma~\ref{lem:dual}, $s_\lambda(h)$ is the size-$L$ determinant
$\det[e_{\lambda'_i-i+j}]$. Every entry has degree at most $1$ in
$\alpha_1$, by~\eqref{eq:e}. To extract the coefficient of $\alpha_1^L$,
replace an entry of index $q\ge1$ by $(-1)^{q-1}\widehat h_{q-1}$,
and an entry of index $q\le0$ by zero. Every determinant term has the
same sign $(-1)^{|\lambda|-L}$.
After removing this sign the determinant is
\[
 \det[\widehat h_{\lambda'_i-1-i+j}]_{1\le i,j\le L}
 =s_{(\lambda'_1-1,\ldots,\lambda'_L-1)}(\widehat h)
 =s_\nu(\widehat h).
\]
An entry of index $q\le0$ is correctly zero because its new moment index
$q-1$ is negative. Trailing zero parts can be removed as noted above.
Finally $|\nu|=|\lambda|-L<|\lambda|$. Induction on partition size,
starting with $s_\varnothing=1$, proves that the extracted coefficient is
nonzero, and therefore the degree bound is exact.
\end{proof}

\section{Frobenius coordinates give the generic leading term}
Apply repeatedly the operation
\[
 \lambda\longmapsto(\lambda_2,\lambda_3,\ldots)'.
\]
At each stage record the first row length. Denote the resulting positive
sequence by $(d_1,\ldots,d_t)$.

\begin{proposition}[Generic Schur--moment leading term]\label{prop:generic}
In lexicographic order $\alpha_1>\alpha_2>\cdots$,
\begin{equation}\label{eq:generic}
 \LT s_\lambda(\mu(\alpha))
 =(-1)^{\sum_{j\ge1}d_{2j}}
   \prod_{h\ge1}\alpha_h^{d_h}.
\end{equation}
The $d_h$ are exactly the sequence in~\eqref{eq:profile}, with a final zero
omitted. In particular, every such generic Schur determinant is nonzero
and has leading coefficient $1$ or $-1$.
\end{proposition}
\begin{proof}
Lemma~\ref{lem:peel} successively determines the greatest possible exponent
of $\alpha_1$, then of $\alpha_2$ among those terms, and so on. It removes
$d_h$ cells at stage $h$, so the total sign exponent is
\[
 \sum_{h\ge1}\sum_{j>h}d_j
 =\sum_{j\ge1}(j-1)d_j
 \equiv\sum_{j\ge1}d_{2j}\pmod2.
\]
For the coordinate formula, after two operations the remaining shape is
$(\lambda_2-1,\lambda_3-1,\ldots)$, with zero parts removed. Thus after
$2i-2$ operations, for $i\le D$, its row lengths are the positive entries of
$(\lambda_i-i+1,\lambda_{i+1}-i+1,\ldots)$; nonpositive entries are omitted.
The next row length is $\lambda_i-i+1=A_i+1$, and the row length after
transposition and removal is $\lambda'_i-i=B_i$. The operations stop
exactly when the shape is empty. This proves both claims.
\end{proof}

\begin{remark}
The proof uses only reciprocal-series algebra and determinant identities.
It does not use positivity or analytic moment theory. The continued
fraction/path framework is classical \cite{Flajolet,Sokal}; the precise
leading-term statement above is proved here, without a priority claim.
\end{remark}

\section{Specialization, signs, and the explicit bijection}
\begin{proof}[Proof of Theorem~\ref{thm:formula}]
Substitute $\ell_{j+1}=y_1\mu_j$ in~\eqref{eq:minor}, transpose the moment
matrix, and reverse both its rows and columns. The two reversal signs
cancel, and~\eqref{eq:lambda} gives
\begin{equation}\label{eq:schurminor}
 \ell_S=(-1)^{r+1}y_1^r s_\lambda(\mu(\alpha)).
\end{equation}
This is valid also when the original Toeplitz minor has structural zeros,
because negative moment indices are defined to be zero.

Under $\alpha_h=\epsilon_h y_hy_{h+1}$ a profile $d$ becomes the exponent
vector
\[
 (r+d_1,d_1+d_2,d_2+d_3,\ldots).
\]
This map on profiles is injective and preserves lexicographic comparisons:
at the first differing profile coordinate the corresponding $y$-exponents
differ in the same direction. Thus distinct monomials do not merge.
Moreover the Schur determinant is homogeneous of degree $|\lambda|$ in
$\alpha$, so all its specialized terms have total degree
$r+2|\lambda|$. Lexicographic and degree-lexicographic leading terms
therefore coincide for this polynomial.

The generic sign from~\eqref{eq:generic} is $(-1)^{\sum_i B_i}$.
The specialized negative weights contribute $(-1)^{\sum_{i\ge2}B_i}$.
Combining these with the prefactor in~\eqref{eq:schurminor}, the final sign
is $(-1)^{r+1+B_1}=1$, because $B_1=r-1$.
This proves~\eqref{eq:leading}. Since $d_1=\lambda_1=\max(S)-r+1$ and
$d_2=r-1$, the first two exponents have the asserted values.
\end{proof}

\begin{proposition}[Image and inverse]\label{prop:bijection}
Formula~\eqref{eq:leading}, together with the empty subset, gives a bijection
onto the set $\mathcal H_n$ in~\eqref{eq:hilbert}.
\end{proposition}
\begin{proof}
For a nonempty partition its positive profile is consecutive, and
\[
 d_i>d_{i+2}\quad\text{whenever }d_i>0,
\]
where missing terms are zero. This follows from the strict decrease of the
Frobenius arm lengths and of the leg lengths. If the leading exponent is
$e=(r+d_1,d_1+d_2,\ldots)$, then $e_1=e_2+1$, and
\[
 e_j-e_{j+1}=d_{j-1}-d_{j+1}>0
\]
through its positive tail, for $j\ge2$. Also $e_2=\max(S)\le n-1$.
Since $e_2\le n-1$, its strictly decreasing positive tail has at most
$n-1$ entries, so the vector has at most $n$ nonzero coordinates.
Thus $e\in\mathcal H_n$.

Here is a constructive inverse. Omit trailing zeros from $e\in\mathcal H_n$.
For $e=(1)$ use $S=\varnothing$. Otherwise put
\begin{equation}\label{eq:inverseprofile}
 d_i=e_{i+1}-e_{i+2}+e_{i+3}-\cdots\quad(1\le i<\text{length}(e)).
\end{equation}
All these numbers are positive, since the tail is strictly decreasing.
Also $d_i-d_{i+2}=e_{i+1}-e_{i+2}>0$.
The two sequences
\[
 A_i=d_{2i-1}-1,\qquad B_i=d_{2i},
\]
with a final $B_i=0$ when needed, are strictly decreasing nonnegative
integer sequences of the same length. They define a unique partition:
for $1\le i\le D$ its row length is $i+A_i$, and for $i>D$ its row
length is $\#\{j:j+B_j\ge i\}$.
To verify this construction, place a diagonal cell at $(i,i)$, its $A_i$
cells to the right, and its $B_i$ cells below, for each $i$. Strict decrease
ensures that the union is a Ferrers diagram and recovers exactly these arm
and leg lengths.

Let $\lambda$ be this partition and $r=\ell(\lambda)=B_1+1=d_2+1$,
where $d_2=0$ if absent. Then set
\begin{equation}\label{eq:inverseS}
 S=\{\lambda_i+r-i:1\le i\le r\}.
\end{equation}
Its elements are distinct positive integers and its largest is
$d_1+r-1=e_1-1\le n-1$. Equations~\eqref{eq:inverseprofile} give
$e_j=d_{j-1}+d_j$ for $j\ge2$, and $r+d_1=e_1$ because
$e_1=e_2+1$. Hence this inverse really recovers $e$ and the original $S$.
\end{proof}

For example, $S=\{2,5,7\}$ gives $\lambda=(5,4,2)$ and
$d=(5,2,3,1)$. Therefore
\[
 \LT(\ell_{\{2,5,7\}})=y_1^8y_2^7y_3^5y_4^4y_5.
\]
Conversely, alternating tails of $(8,7,5,4,1)$ recover $(5,2,3,1)$.
The rank $r$ can also be recovered directly as $e_1-e_2+e_3-e_4+\cdots$.

\section{The Hilbert basis and an equal-growth principle}
\begin{lemma}[Elementary Hilbert-basis description]\label{lem:hilbert}
The indecomposable nonzero elements of $L_n$ are precisely $\mathcal H_n$.
Every $x\in L_n$ is a sum of elements of $\mathcal H_n$.
\end{lemma}
\begin{proof}
For $n=1$ the assertion is immediate. Suppose $n\ge2$ and put
$g=x_1-x_2$. The first normalized inequality implies
$x_2\le(n-1)g$. Hence $g\ge1$ if $x\ne0$.
The integer points with $g=1$ are exactly~\eqref{eq:hilbert}:
$x_2\le n-1$, and any two consecutive positive coordinates must decrease
strictly. Conversely, an integer vector of the displayed form belongs to
$C_n$. Indeed $x_{i+1}\le n-i$, and whenever it is positive,
$x_i-x_{i+1}\ge1$; therefore
$(n-i)(x_i-x_{i+1})\ge x_{i+1}$, which is the normalized inequality.

For a general nonzero $x$ construct $g$ vectors, indexed by $t=0,\ldots,g-1$.
For $i\ge2$ put
\[
 b_i^{(t)}=\min\{n-i+1,\max\{0,x_i-t(n-i+1)\}\},\qquad
 b_1^{(t)}=b_2^{(t)}+1.
\]
The normalized tails equal
$\min\{1,\max\{0,x_i/(n-i+1)-t\}\}$, so they remain weakly decreasing.
The first inequality holds since $b_2^{(t)}\le n-1$.
Thus $b^{(t)}\in L_n$ and its first-coordinate difference is $1$; it is
an element of $\mathcal H_n$.
For $i\ge2$, the original inequalities give
$0\le x_i\le g(n-i+1)$, and the elementary truncation identity gives
$\sum_t b_i^{(t)}=x_i$. The first coordinates sum to $x_2+g=x_1$.
This proves generation. Finally a vector with difference $1$ cannot split
into two nonzero vectors, since each has positive integral difference.
\end{proof}

The following principle is useful beyond the present problem.

\begin{lemma}[Saturation from equal leading growth]\label{lem:growth}
Let $C\subseteq\mathbb R_{\ge0}^n$ be a full-dimensional rational polyhedral
cone and $L=C\cap\Z^n$. Let $A\subseteq k[y_1,\ldots,y_n]$ be a subalgebra,
with degree-lexicographic initial algebra. Suppose
\[
 k[L]\subseteq\inA A,
\]
and, for the filtration by total degree at most $N$,
\[
 \dim_k A_{\le N}=cN^n+O(N^{n-1}),\qquad
 \#\{e\in L:|e|\le N\}=cN^n+O(N^{n-1}).
\]
Then $\inA A=k[L]$.
\end{lemma}
\begin{proof}
Gaussian elimination on the finite-dimensional vector space $A_{\le N}$
gives a basis with distinct leading monomials. Consequently its dimension
is the number of monomials of degree at most $N$ in $\inA A$.
For the reverse direction in this equality, a product of leading monomials
is the leading monomial of the product of the underlying polynomials;
degree-lexicographic order ensures the latter polynomial has the same
total degree as that leading monomial.

Suppose that $y^u\in\inA A$ with $u\notin C$. Choose one of the defining
linear inequalities $f\ge0$ for $C$ that is violated by $u$, so $f(u)<0$.
Choose linearly independent integer vectors $v_1,\ldots,v_n\in C$.
Fix small positive constants $\kappa,\eta$ such that
\[
 \kappa|u|+\eta\sum_i|v_i|<1,
 \qquad \kappa f(u)+\eta\sum_i f(v_i)<0.
\]
This is possible by first choosing $\kappa$ and then $\eta$.
For all sufficiently large $N$, the vectors
\[
 \lfloor\kappa N\rfloor u+\sum_{i=1}^n m_iv_i,
 \qquad 0\le m_i\le\lfloor\eta N\rfloor,
\]
have total degree at most $N$, violate $f\ge0$, and belong to the exponent
semigroup of $\inA A$. They are distinct because the $v_i$ are independent.
Thus, in addition to all the lattice points in $C$ of degree at most $N$,
the initial algebra contains at least
$(\lfloor\eta N\rfloor+1)^n$ other such monomials. This contradicts the
equal leading growth constants. Hence no $u$ outside $C$ occurs. Since
$L$ contains every lattice point in $C$, the asserted equality follows.
\end{proof}

\begin{lemma}[The two growth constants]\label{lem:volumes}
For the Lecture Hall algebra and cone,
\begin{align*}
 \dim_k A_n(k)_{\le N}
 &=\frac{2^n}{(2n)!}N^n+O(N^{n-1}),\\
 \#\{e\in L_n:|e|\le N\}
 &=\frac{2^n}{(2n)!}N^n+O(N^{n-1}).
\end{align*}
\end{lemma}
\begin{proof}
The unique Dyck path of semilength $i-1$ that reaches height $i-1$ is the
single mountain. Thus for $i\ge2$
\begin{equation}\label{eq:triangular}
 \ell_i=c_i y_i+q_i(y_1,\ldots,y_{i-1}),\qquad
 c_i=\left(\prod_{h=1}^{i-1}\epsilon_h\right)
          \prod_{j=1}^{i-1}y_j^2\ne0.
\end{equation}
Together with $\ell_1=y_1$, this proves algebraic independence over every
field: in a putative relation viewed as a polynomial in $\ell_i$, the
coefficient of the highest power of $y_i$ is nonzero by induction.
Each $\ell_i$ has total degree $2i-1$. Therefore the first count equals
\[
 \#\{a\in\N^n:\sum_{i=1}^n(2i-1)a_i\le N\}
 =\frac{N^n}{n!\prod_{i=1}^n(2i-1)}+O(N^{n-1})
 =\frac{2^n}{(2n)!}N^n+O(N^{n-1}).
\]
The estimate counts integer points in a weighted simplex. It follows, for
example, by attaching a unit cube to each nonnegative integer point:
the union lies between the weighted simplices of bounds $N$ and
$N+\sum_i(2i-1)$.

For the second count the cone is generated by the independent integer rays
\[
 v^{(j)}=(n,n-1,\ldots,n-j+1,0,\ldots,0),\qquad 1\le j\le n.
\]
If $V$ is their column matrix, $|\det V|=n!$, and
$w_j=|v^{(j)}|=j(2n-j+1)/2$. Hence
\[
 \vol\{x\in C_n:|x|\le1\}
 =\frac{|\det V|}{n!\prod_j w_j}
 =\frac{2^n}{(2n)!},
 \qquad \prod_jw_j=\frac{(2n)!}{2^n}.
\]
Integer-point count equals this volume times $N^n$, up to $O(N^{n-1})$.
Here is an elementary justification: cover the dilated simplex by unit
lattice cubes. Cubes wholly inside contribute their full volume and one
lattice point under a fixed half-open convention. Only cubes intersecting
the boundary can affect the difference. Each of its finitely many facets
intersects $O(N^{n-1})$ cubes: solve its hyperplane equation for a coordinate
with nonzero coefficient after fixing the other $n-1$ cube coordinates.
There are $O(N^{n-1})$ such choices and a bounded number of remaining
coordinate choices. This proves the required error estimate.
\end{proof}

\section{The initial algebra and the Lecture Hall series}
\begin{proof}[Proof of Theorem~\ref{thm:main}]
Every $\ell_S$ belongs to $A_n(k)$ by its determinant definition.
Theorem~\ref{thm:formula} and Proposition~\ref{prop:bijection} show that
its monic leading terms run through $y^h$, $h\in\mathcal H_n$.
Lemma~\ref{lem:hilbert} therefore gives
$k[L_n]\subseteq\inA A_n(k)$.
Lemmas~\ref{lem:volumes} and~\ref{lem:growth} prove the reverse inclusion.
The leading monomials of $\mathcal B_n$ generate this initial algebra, so
$\mathcal B_n$ is a SAGBI basis by definition. Since the exponent vectors
are distinct indecomposable elements of the semigroup, none can be
removed. This also proves minimality.
\end{proof}

\begin{corollary}[Bivariate Lecture Hall Theorem]\label{cor:series}
As a formal power series,
\[
 \sum_{e\in L_n}q_1^{e_1+e_3+\cdots}q_2^{e_2+e_4+\cdots}
 =\prod_{i=1}^n\frac{1}{1-q_1^i q_2^{i-1}}.
\]
\end{corollary}
\begin{proof}
Give odd-indexed $y_j$ bidegree $(1,0)$ and even-indexed $y_j$ bidegree
$(0,1)$. Every adjacent weight $y_hy_{h+1}$ has bidegree $(1,1)$, so
$\ell_i$ has bidegree $(i,i-1)$. Algebraic independence gives the product
on the right as the Hilbert series of $A_n(k)$. Gaussian elimination
within each finite-dimensional bihomogeneous component shows that taking
initial terms preserves this Hilbert series. Theorem~\ref{thm:main}
identifies the initial algebra with the monomial algebra on the left.
\end{proof}

This is the classical theorem of Bousquet--M\'elou and Eriksson, recalled
in \cite[Theorem 1.2]{Katthan}; the identity itself is not new. The proof
here derives it from the initial-algebra result and elementary volume
comparison, rather than using it to prove that result. In particular,
the argument has no circular dependence on the Lecture Hall Theorem.

\section{An integral basis and an explicit flat family}
Set $A_n(\Z)=\Z[\ell_1,\ldots,\ell_n]$.

\begin{theorem}[Integral and deformation extensions]\label{thm:flat}
The set $\mathcal B_n$ is a monic SAGBI basis over $\Z$ and
\[
 \inA A_n(\Z)=\Z[L_n].
\]
There is an explicit finitely generated $\Z[t]$-subalgebra
$\mathcal R_n\subseteq\Z[t,y_1,\ldots,y_n]$, free as a $\Z[t]$-module,
with
\[
 \mathcal R_n/(t-1)\mathcal R_n\cong A_n(\Z),\qquad
 \mathcal R_n/t\mathcal R_n\cong\Z[L_n].
\]
The same construction remains free, hence flat, after extension to any
field. Thus it supplies the toric deformation in all characteristics.
\end{theorem}
\begin{proof}
Embed the integral algebra in $A_n(\Q)$. The leading exponent of any
nonzero integral polynomial in the algebra belongs to $L_n$ by
Theorem~\ref{thm:main}. For each $e\in L_n$, choose a decomposition
$e=h_1+\cdots+h_a$ into elements of $\mathcal H_n$, for example the one
in Lemma~\ref{lem:hilbert}. Use Proposition~\ref{prop:bijection} to select
$S(h_j)$ and put
\[
 F_e=\prod_{j=1}^a\ell_{S(h_j)},\qquad F_0=1.
\]
It has leading monomial $y^e$ with coefficient $1$.
Repeatedly subtracting an integer multiple of the corresponding $F_e$
from a polynomial cancels its leading term and strictly decreases the
leading monomial. The process terminates: total degrees do not increase,
and there are finitely many monomials in a fixed degree bound. It uses no
division. Thus the $F_e$ span $A_n(\Z)$; their distinct leading monomials
also make them linearly independent over $\Z$.
This proves the integral SAGBI assertion.

For an explicit weight vector, let $D_n=n(n+1)/2$ and set
\[
 W_j=(D_n+1)^{n-j}\qquad(1\le j\le n).
\]
Every generator $\ell_S$ has total degree at most $D_n$.
Indeed for $r=|S|>0$ the degree is
$2\sum_{s\in S}s+2r-r^2\le r(2n+1-2r)\le D_n$;
the last inequality follows because
$D_n-r(2n+1-2r)=(n-2r)(n-2r+1)/2\ge0$.
The empty subset has degree $1$.
These weights make the lexicographic leading monomial of each generator
its unique greatest-weight monomial. At the first differing exponent
coordinate, its positive contribution is at least $W_j$, whereas all
possible negative contributions from later coordinates are at most
$D_nW_{j+1}<W_j$.

For $h\in\mathcal H_n$ define
\[
 \widetilde\ell_h(t,y)=t^{W\cdot h}
      \ell_{S(h)}(t^{-W_1}y_1,\ldots,t^{-W_n}y_n),
 \quad
 \mathcal R_n=\Z[t,\widetilde\ell_h:h\in\mathcal H_n].
\]
The displayed expressions are polynomials: every occurring monomial has
weight at most $W\cdot h$. For each $e\in L_n$ similarly put
\[
 \widetilde F_e=t^{W\cdot e}F_e(t^{-W_1}y_1,\ldots,t^{-W_n}y_n).
\]
It is a product of the indicated generators.

Each $F_e$ has unique highest-weight monomial $y^e$. In the unique integral
basis expansion of a polynomial, the largest weight among the occurring
$F_e$ cannot cancel: their highest-weight monomials are distinct. Since
$F_eF_f-F_{e+f}$ has weight strictly below $W\cdot(e+f)$, it follows that
\[
 F_eF_f=F_{e+f}+\sum_{W\cdot g<W\cdot(e+f)}c_gF_g,
 \qquad c_g\in\Z.
\]
Homogenizing yields
\[
 \widetilde F_e\widetilde F_f
 =\widetilde F_{e+f}+
   \sum_g c_g t^{W\cdot(e+f)-W\cdot g}\widetilde F_g.
\]
Thus the $\Z[t]$-span of the $\widetilde F_e$ is closed under multiplication
and is exactly $\mathcal R_n$. Their distinct leading $y$-monomials, each
with coefficient $1$, show that this span is free over $\Z[t]$.
Modulo $t$ the displayed multiplication becomes
$\widetilde F_e\widetilde F_f=\widetilde F_{e+f}$, giving $\Z[L_n]$.
At $t=1$ the basis specializes to the basis $F_e$ of $A_n(\Z)$.
After inverting $t$, the invertible change of variables
$y_j\mapsto t^{-W_j}y_j$ and invertible scalar factors on the generators give
$\mathcal R_n[t^{-1}]\cong A_n(\Z)[t,t^{-1}]$.
Base change preserves this explicit free basis and both fiber descriptions.
\end{proof}

\section{Finite records and scope}\label{sec:verification}
The accompanying data contain finite exact records with the scopes below.
The full extension suite and the earlier polynomiality suite were not rerun
for this edition. The recorded regeneration of the extension data was
completed in separate stages after an interruption; original process logs
are not included. Stored counts and hashes do not independently certify
that execution history.

The verifier reconstructs the finite tests summarized below.
Its input is the Dyck-path definition and the original Toeplitz minors;
the leading-term formula is used only as the expected answer.
\begin{center}
\begin{tabular}{p{0.36\textwidth}p{0.55\textwidth}}
\toprule
Check & Exact scope and output\\
\midrule
Original flag minors & All $2^{10}=1024$ subsets of $[10]$, including the
empty subset, in $11$ variables. All leading exponents and coefficients
agree. The expanded polynomials have $884566$ terms in total; the largest
has $5427$ terms.\\[3pt]
Generic Schur determinants & All $195$ partitions of sizes $0$ through $11$,
computed by exact Bareiss determinants of independently generated moments.\\[3pt]
Bijection and inverse & Every subset for each $1\le n\le16$: $65535$ formula
checks. These check the combinatorial formulas, not symbolic determinants
in dimension $16$.\\[3pt]
Lattice and Hilbert counts & Direct enumeration in each $1\le n\le8$ through
total degree $40$: $24485$ vectors overall. Bigraded counts agree with the
weighted polynomial algebra, and each layer decomposition is checked.\\[3pt]
Full small polynomials & All $128$ minors for subsets of $[7]$ are recorded
as complete sparse coefficient lists in $8$ variables.\\
\bottomrule
\end{tabular}
\end{center}

The directory \texttt{code/} contains the executable verifier and comparison
script. The directory \texttt{data/} contains all expected outputs, including
leading exponents, partition profiles, term counts, and canonical polynomial
hashes for the larger symbolic tests. Hash formatting is documented in the
top-level README. Running the verifier regenerates the six mathematical
data files. The earlier polynomiality manuscript and its separate
verification files remain at the top level of this research package.

The proofs, not the finite tests, establish the all-$n$ statements.
The work does not claim that a numerical specialization retains algebraic
independence or a SAGBI basis, and does not address uniqueness of the full
defining system at vanishing divisors. Its algebra statements concern
independent indeterminates. The generic Schur lemma, the explicit bijection,
and the all-dimensional conclusions are proved above; they are not listed
as conjectures or as extrapolations. A limited literature check located the
original question and its publication record but does not certify priority
for these arguments. Independent mathematical review remains appropriate.

\paragraph{AI assistance.}
This extension was developed with AI assistance. All mathematical reasoning
needed for the stated results is presented here, and the exact verification
code and data are provided for independent review.

\appendix
\section{Identification with the defining Toeplitz equations}\label{app:identification}
This appendix reproduces the foundation established in \cite{Polynomiality}.
It is included for completeness, not claimed as a new extension.
For a sequence $P$, let
\[
 M(P)_{r,c}=\begin{cases}-P_{c-r+1}&c\ge r,\\0&c<r.\end{cases}
\]
Define the square minors and target monomials by
\[
 E_i(P)=-\det M(P)_{\{1,\ldots,\lceil i/2\rceil\},
                       \{\lfloor i/2\rfloor+1,\ldots,i\}},\qquad
 T_i=\prod_{j=1}^i y_j^{i+1-j}.
\]

\begin{lemma}[Hankel products]\label{lem:hankel}
For independent Dyck weights and every $m\ge1$,
\begin{align*}
 H_m^{(0)}:=\det[\mu_{i+j}]_{0\le i,j<m}
 &=\prod_{h=1}^{m-1}(\alpha_{2h-1}\alpha_{2h})^{m-h},\\
 H_m^{(1)}:=\det[\mu_{i+j+1}]_{0\le i,j<m}
 &=\alpha_1^m\prod_{h=1}^{m-1}(\alpha_{2h}\alpha_{2h+1})^{m-h}.
\end{align*}
\end{lemma}
\begin{proof}
Pair consecutive steps of a Dyck path. This gives a Motzkin path with
up weight $1$, down weight $\lambda_h=\alpha_{2h-1}\alpha_{2h}$,
and horizontal weights $b_0=\alpha_1$ and
$b_h=\alpha_{2h}+\alpha_{2h+1}$ for $h\ge1$.
Let $J$ be the tridiagonal matrix on nonnegative indices, with diagonal
$b_h$, superdiagonal $1$, and subdiagonal $\lambda_h$.
Then $\mu_j=(J^j)_{0,0}$ by path enumeration.
Put $d_0^*=1$, $d_h^*=\lambda_1\cdots\lambda_h$, and
$D=\diag(d_0^*,d_1^*,\ldots)$. Direct multiplication gives
$JD=DJ^\transpose$, hence
$(J^j)_{h,0}=d_h^*(J^j)_{0,h}$.

Fix $m$ and write $L_{ih}=(J^i)_{0,h}$ for $0\le i,h<m$.
This is unit lower triangular. If $J_m,D_m$ are the leading blocks, path
length bounds give the exact factorizations
\[
 [\mu_{i+j}]_{i,j<m}=LD_mL^\transpose,\qquad
 [\mu_{i+j+1}]_{i,j<m}=LJ_mD_mL^\transpose.
\]
For the second identity split a path after $i$ and then one more step;
the two interface heights are at most $i$ and $j$, respectively, so no
term outside these finite blocks occurs. The first determinant is now
$\prod_{h=0}^{m-1}d_h^*$.

For $\delta_m=\det J_m$, tridiagonal expansion gives
\[
 \delta_0=1,\quad\delta_1=\alpha_1,\quad
 \delta_m=(\alpha_{2m-2}+\alpha_{2m-1})\delta_{m-1}
          -\alpha_{2m-3}\alpha_{2m-2}\delta_{m-2}.
\]
Induction cancels the two terms involving $\alpha_{2m-2}$ and yields
$\delta_m=\prod_{h=1}^m\alpha_{2h-1}$.
Taking the second determinant and rearranging factors gives the shifted
product. No division or sign assumption is used.
\end{proof}

\begin{proposition}[The moment identification]\label{prop:identification}
The polynomials in~\eqref{eq:weights} satisfy $E_i(\ell)=T_i$ for every
$i\ge1$. Over the fraction field of the independent $y_j$, they are the
unique sequence satisfying those equations.
\end{proposition}
\begin{proof}
Put $s_m=(-1)^{1+m(m+1)/2}$. Row reversal, row signs, and the outer minus
in the Toeplitz minors give
\[
 E_{2m-1}(\ell)=s_m y_1^mH_m^{(0)},\qquad
 E_{2m}(\ell)=s_m y_1^mH_m^{(1)}.
\]
Substitute~\eqref{eq:weights} in Lemma~\ref{lem:hankel}.
The exponents of $y_j$ in the two products, after multiplication by $y_1^m$,
are respectively $2m-j$ and $2m+1-j$ on their supports. Their common sign,
for $m\ge2$, is
\[
 (-1)^{\sum_{h=2}^{m-1}(m-h)}
 =(-1)^{(m-2)(m-1)/2}=s_m.
\]
For $m=1$ both signs are $1$ directly. This proves all the equations.

For uniqueness, $E_1=P_1$ and $E_2=P_2$. For $i\ge3$ the entry $-P_i$
occurs only in the first row and last column of its square minor. If
$k=\lceil i/2\rceil$, its coefficient in $E_i$ is
$(-1)^kE_{i-2}(P)$, by cofactor expansion and shifting indices.
After the previous equations have been solved this coefficient is
$(-1)^kT_{i-2}\ne0$. Successive solution in the fraction field gives a
unique new $P_i$ each time. This identifies the sequence with the one in
the cited polynomiality note and source. The argument does not assert stepwise uniqueness
after a specialization that makes these coefficients vanish.
\end{proof}

\begin{thebibliography}{9}
\bibitem{Katthan}
Lukas Katth\"an,
\emph{The Lecture Hall Cone as a toric deformation},
arXiv:1809.01377v1 (2018), especially Definitions 3.1 and 4.1,
Conjecture 4.2, and Problem 5.1.
\url{https://arxiv.org/abs/1809.01377v1}.
Published version: \emph{Algebraic and Geometric Combinatorics on Lattice
Polytopes}, World Scientific (2019), 242--253,
\url{https://doi.org/10.1142/9789811200489_0015}.
The theorem numbering used throughout this report is the arXiv numbering.

\bibitem{Polynomiality}
\emph{Polynomiality of the Lecture Hall sequence via signed Dyck paths},
version 0.1.1, October 13, 2026.
\url{https://github.com/Aurantium-47/lecture-hall-polynomiality/releases/tag/v0.1.1}.
The moment identification is reproduced in Appendix~\ref{app:identification}.

\bibitem{Sokal}
Alan D. Sokal,
\emph{A simple algorithm for expanding a power series as a continued fraction},
Expositiones Mathematicae \textbf{41} (2023), 245--287.
\url{https://arxiv.org/abs/2206.15434}.
\url{https://doi.org/10.1016/j.exmath.2022.12.001}.

\bibitem{Flajolet}
Philippe Flajolet,
\emph{Combinatorial aspects of continued fractions},
Discrete Mathematics \textbf{32} (1980), 125--161.
\url{https://algo.inria.fr/flajolet/Publications/Flajolet80b.pdf}.
\end{thebibliography}
\end{document}
